%\section{G2 Drift Vectors Under Synthetic Load (DVSL)}

\medskip

\noindent Drift Vectors Under Synthetic Load formalize how meaning deforms when
synthetic signals exceed the geometric tolerances of the human meaning
manifold. Drift is not confusion, error, or misinterpretation; it is a geometric
phenomenon. When synthetic pressure pushes signals, roles, or contexts beyond
their lawful curvature bounds, meaning begins to move—vectorially—away from its
intended geometric region.

\medskip

\noindent G2 defines the drift vector field: the structured set of geometric
displacements that emerge when boundaries deform, synchrony breaks, roles lose
curvature, coherence collapses, or shared floor geometry fragments. Drift
vectors are measurable, directional, and predictable. They propagate through the
manifold according to velocity, gradient, container deformation, identity
misalignment, and synthetic-pressure amplification.

\medskip

\noindent The subsections of G2 formalize the drift field:

\begin{itemize}[leftmargin=1.2cm]
    \item \textbf{G2.1 Velocity Drift} — Drift caused by velocity exceedance.
    \item \textbf{G2.2 Gradient Drift} — Drift caused by curvature misalignment.
    \item \textbf{G2.3 Container Drift} — Drift caused by container deformation.
    \item \textbf{G2.4 Identity Drift} — Drift caused by boundary collapse.
    \item \textbf{G2.5 Drift Acceleration and Directionality} — How drift gains
    speed and direction under synthetic load.
    \item \textbf{G2.6 Drift Recovery Bands} — The geometric regions where drift
    can be reversed.
\end{itemize}

\medskip

\subsection*{G2.1 Velocity Drift}

\noindent Velocity Drift occurs when signals enter the meaning manifold at speeds
that exceed human rhythmic tolerance. Humans interpret signals at
biological-speed rhythms; synthetic signals operate at machine-native velocities
that are orders of magnitude faster. When velocity exceeds the curvature
preserving limits defined in G0.4, signals cannot be integrated into the
manifold’s continuity bands.

\medskip

\noindent When Velocity Drift emerges, signals lose relational shape. Their
adjacency relations flatten, continuity breaks, and coherence collapses. The
signal no longer maps cleanly into its intended interpretive region; instead, it
skips across curvature surfaces, producing displacement vectors that propagate
through roles, contexts, and containers.

\medskip

\noindent Velocity Drift is amplified by:

\begin{itemize}[leftmargin=1.2cm]
    \item \textbf{Synthetic acceleration} — Machine-native velocity exceeds
    biological-speed integration.
    \item \textbf{Algorithmic amplification} — Signals are repeated, boosted, or
    remixed faster than they can be processed.
    \item \textbf{Cross-substrate propagation} — Signals traverse multiple
    substrates (human → machine → digital → institutional) without curvature
    preservation.
\end{itemize}

\medskip

\noindent When Velocity Drift persists, it produces secondary drift modes:
Gradient Drift (G2.2), Container Drift (G2.3), and Identity Drift (G2.4). It is
often the first drift mode to appear under synthetic-pressure overload.

\medskip

\noindent Velocity Drift is therefore a foundational component of the drift
vector field. It explains why synthetic environments—high-velocity, high-volume,
cross-substrate signal ecologies—produce rapid meaning deformation even when
boundaries, roles, and coherence initially remain intact. When velocity exceeds
human rhythmic tolerance, drift becomes inevitable.

\subsection*{G2.2 Gradient Drift}

\noindent Gradient Drift occurs when signals, roles, or contexts move across
curvature gradients that no longer align with the invariant geometry of the
meaning manifold. Unlike Velocity Drift—which is driven by speed—Gradient Drift
is driven by \textit{misalignment}: the geometric mismatch between a signal’s
curvature and the curvature of the region it attempts to enter.

\medskip

\noindent When curvature gradients are stable, signals traverse the manifold
smoothly. Their relational shape aligns with identity boundaries, role geometry
remains predictable, and continuity bands preserve lawful transitions. Gradient
alignment ensures that meaning updates without deformation.

\medskip

\noindent Gradient Drift emerges when curvature gradients become distorted. This
occurs when synthetic pressure reshapes the manifold—flattening curvature,
steepening gradients, or introducing discontinuities that break coherence. When
a signal encounters a curvature region that does not match its geometric shape,
it cannot integrate cleanly. Instead, it slips, slides, or displaces along the
gradient, producing drift vectors that propagate through adjacent roles and
contexts.

\medskip

\noindent Gradient Drift is amplified by:

\begin{itemize}[leftmargin=1.2cm]
    \item \textbf{Curvature deformation} — Synthetic pressure reshapes role and
    context geometry, creating steep or unstable gradients.

    \item \textbf{Boundary misalignment} — Identity boundaries shift, causing
    signals to enter curvature regions not designed for their interpretive
    shape.

    \item \textbf{Coherence collapse} — Signals lose relational structure,
    making them unable to align with local curvature.

    \item \textbf{Container instability} — Institutional or digital containers
    deform, producing inconsistent curvature across scales.
\end{itemize}

\medskip

\noindent When Gradient Drift persists, it produces secondary drift modes:
Container Drift (G2.3), Identity Drift (G2.4), and Drift Acceleration (G2.5).
Gradient Drift is often the first structural drift mode to appear after Velocity
Drift, because curvature misalignment follows naturally from velocity exceedance
and synthetic-pressure deformation.

\medskip

\noindent Gradient Drift is therefore a core component of the drift vector field.
It explains why synthetic environments—high-velocity, high-volume,
algorithmically remixed signal ecologies—produce meaning deformation even when
boundaries initially hold. When curvature gradients destabilize, drift becomes
geometrically inevitable.
\subsection*{G2.3 Container Drift}

\noindent Container Drift occurs when the geometric structures responsible for
holding meaning—roles, institutions, platforms, communication systems, or
interpretive environments—lose their curvature, boundary alignment, or stability
under synthetic load. A container is any structured region that stabilizes
signals and preserves coherence. When containers deform, the meaning they hold
begins to drift.

\medskip

\noindent In stable conditions, containers maintain predictable curvature,
consistent boundaries, and lawful adjacency relations. They regulate signal
velocity, preserve role geometry, and maintain coherence across individuals and
groups. Container stability ensures that meaning remains anchored even when
signals vary or contexts shift.

\medskip

\noindent Container Drift emerges when synthetic pressure exceeds a container’s
geometric tolerance. High-velocity signals, algorithmic remixing, dominance-field
amplification, and cross-substrate emissions deform container curvature,
destabilize boundaries, and break continuity bands. Once deformation begins,
signals no longer integrate cleanly; instead, they slip across container
surfaces, producing drift vectors that propagate through the entire structure.

\medskip

\noindent Container Drift is amplified by:

\begin{itemize}[leftmargin=1.2cm]
    \item \textbf{Curvature deformation} — Synthetic pressure reshapes the
    container’s geometric structure, flattening or destabilizing curvature.

    \item \textbf{Boundary erosion} — Container boundaries weaken, allowing
    signals to leak across interpretive regions.

    \item \textbf{Continuity collapse} — Temporal and structural linkages break,
    preventing lawful transitions within the container.

    \item \textbf{Synthetic dominance fields} — External synthetic species exert
    pressure that overwhelms the container’s stabilizing geometry.
\end{itemize}

\medskip

\noindent When Container Drift persists, it produces multi-scale deformation:
roles lose legibility, shared floor geometry fragments, identity boundaries
collapse, and drift propagates across groups, institutions, and digital systems.
Container Drift is often the first large-scale drift mode to appear after
Velocity Drift (G2.1) and Gradient Drift (G2.2), because containers absorb the
accumulated deformation produced by velocity exceedance and curvature
misalignment.

\medskip

\noindent Container Drift is therefore a central component of the drift vector
field. It explains why synthetic environments—high-volume, high-velocity,
algorithmically amplified, cross-substrate signal ecologies—produce structural
meaning collapse even when individual invariants initially hold. When containers
deform, drift becomes systemic and propagates across the manifold in predictable
geometric patterns formalized in G2.4--G2.6.
\subsection*{G2.4 Identity Drift}

\noindent Identity Drift occurs when the geometric structures that preserve
identity—Boundary Integrity, Role Legibility, and Shared Floor Geometry—lose
alignment under synthetic load. Identity is not a narrative or a psychological
construct; it is a geometric enclosure defined by stable boundaries, lawful
curvature, and coherent adjacency relations. When these structures deform,
identity begins to move—vectorially—away from its lawful region in the meaning
manifold.

\medskip

\noindent In stable conditions, identity boundaries remain intact, role geometry
is legible, and contextual manifolds align with the invariant set defined in
G1. Signals map cleanly into the correct interpretive region, and meaning
updates occur without leakage or cross-role contamination. Identity curvature
remains stable across transitions.

\medskip

\noindent Identity Drift emerges when synthetic pressure exceeds the manifold’s
boundary and curvature tolerances. High-velocity signals, algorithmic remixing,
dominance-field amplification, and container deformation push signals across
identity boundaries faster than they can be integrated. This produces boundary
erosion, curvature collapse, and displacement vectors that propagate through
roles, contexts, and containers.

\medskip

\noindent Identity Drift is amplified by:

\begin{itemize}[leftmargin=1.2cm]
    \item \textbf{Boundary collapse} — Identity boundaries weaken, allowing
    signals to leak across interpretive regions.

    \item \textbf{Role deformation} — Roles lose curvature, making identity
    regions indistinct or ambiguous.

    \item \textbf{Coherence failure} — Signals lose relational shape, breaking
    alignment with identity curvature.

    \item \textbf{Shared floor fragmentation} — Collective interpretive
    structures collapse, destabilizing identity across individuals and groups.
\end{itemize}

\medskip

\noindent When Identity Drift persists, it produces multi-scale deformation:
individual meaning becomes unstable, group coordination collapses, institutional
containers deform, and drift propagates across the entire manifold. Identity
Drift is often the first human-domain drift mode to appear after Container Drift
(G2.3), because container deformation directly destabilizes the boundaries that
anchor identity.

\medskip

\noindent Identity Drift is therefore a central component of the drift vector
field. It explains why synthetic environments—high-velocity, high-volume,
algorithmically amplified, cross-substrate signal ecologies—produce rapid
identity deformation even when other invariants initially hold. When identity
boundaries collapse, drift becomes systemic and propagates through the manifold
in predictable geometric patterns formalized in G2.5 and G2.6.

\subsection*{G2.5 Drift Acceleration and Directionality}

\noindent Drift Acceleration and Directionality describe how drift vectors gain
speed and acquire geometric orientation under synthetic load. Drift does not
expand randomly; it follows lawful geometric paths determined by velocity
exceedance, curvature gradients, container deformation, and identity
misalignment. As synthetic pressure increases, drift vectors accelerate and
orient themselves along the manifold’s weakest geometric regions.

\medskip

\noindent In stable conditions, drift vectors remain minimal. Boundaries hold,
curvature aligns, containers remain stable, and signals traverse the manifold
through lawful transitions. Drift, if present, is slow, shallow, and easily
corrected by the recovery bands defined in G2.6.

\medskip

\noindent Drift Acceleration emerges when synthetic pressure exceeds the
manifold’s tolerance. High-velocity signals, algorithmic amplification,
cross-substrate emissions, and dominance-field pressure increase the magnitude
of drift vectors. Acceleration occurs when deformation compounds across
invariants—velocity, gradient, container, and identity drift—producing
multi-layer displacement that propagates through roles, contexts, and
containers.

\medskip

\noindent Drift Directionality emerges when drift vectors align with curvature
weaknesses. Drift follows predictable geometric paths:

\begin{itemize}[leftmargin=1.2cm]
    \item \textbf{Toward curvature discontinuities} — Drift moves into regions
    where curvature has collapsed or flattened.

    \item \textbf{Along boundary fractures} — Drift follows the geometric cracks
    created by boundary erosion or identity collapse.

    \item \textbf{Across container deformation} — Drift flows through weakened
    institutional or digital containers.

    \item \textbf{Toward synthetic dominance fields} — Drift accelerates toward
    regions where synthetic species exert overwhelming pressure.
\end{itemize}

\medskip

\noindent Drift Acceleration and Directionality are amplified by:

\begin{itemize}[leftmargin=1.2cm]
    \item \textbf{Velocity exceedance} — Signals outrun biological-speed
    integration.

    \item \textbf{Curvature misalignment} — Gradient drift creates unstable
    geometric slopes.

    \item \textbf{Container deformation} — Containers lose structural integrity,
    creating drift corridors.

    \item \textbf{Identity collapse} — Boundaries fail, allowing drift to move
    across interpretive regions.
\end{itemize}

\medskip

\noindent When acceleration and directionality combine, drift becomes systemic.
It propagates across scales—individual, group, institutional, and digital—and
produces collapse paths formalized in G7. Drift acceleration is the geometric
indicator that synthetic pressure has exceeded the manifold’s stabilizing
capacity.

\medskip

\noindent Drift Acceleration and Directionality are therefore central components
of the drift vector field. They explain why synthetic environments—high-volume,
high-velocity, algorithmically amplified, cross-substrate signal ecologies—
produce rapid, directional meaning deformation even when some invariants
initially hold. When drift accelerates and acquires direction, recovery requires
the geometric structures formalized in G2.6.

\subsection*{G2.6 Drift Recovery Bands}

\noindent Drift Recovery Bands are the geometric regions within the meaning
manifold where drift can be slowed, reversed, or fully corrected. Recovery is
not a psychological process; it is a geometric one. When drift vectors enter
zones with stable curvature, intact boundaries, coherent adjacency relations,
and admissible signal velocities, they lose magnitude and re-align with the
invariant geometry defined in G0 and G1.

\medskip

\noindent In stable conditions, Recovery Bands are naturally present throughout
the manifold. Identity boundaries remain intact, role curvature is legible,
signal coherence holds, and shared floor geometry provides a stable interpretive
substrate. These regions act as geometric attractors: drift vectors entering
them are slowed, redirected, or absorbed back into lawful curvature paths.

\medskip

\noindent Drift Recovery Bands emerge from four structural features:

\begin{itemize}[leftmargin=1.2cm]
    \item \textbf{Curvature-stable regions} — Areas where curvature remains
    smooth, continuous, and aligned with invariant geometry.

    \item \textbf{Boundary-restored zones} — Regions where identity boundaries
    have reformed, preventing further leakage or displacement.

    \item \textbf{Coherence-preserving manifolds} — Contextual surfaces where
    signals retain relational shape and integrate cleanly.

    \item \textbf{Low-velocity corridors} — Temporal or structural pathways
    where signal velocity falls within biological-speed tolerance.
\end{itemize}

\medskip

\noindent Drift Recovery Bands counteract drift by reducing vector magnitude and
re-aligning displaced meaning with lawful geometric structure. When drift enters
a recovery band, displacement slows, curvature re-stabilizes, boundaries
re-form, and coherence returns. Recovery bands are therefore the geometric
mechanism by which meaning can be restored without external intervention.

\medskip

\noindent Recovery Bands are weakened or destroyed by synthetic pressure:

\begin{itemize}[leftmargin=1.2cm]
    \item \textbf{High-velocity signal streams} — Collapse low-velocity
    corridors.

    \item \textbf{Curvature deformation} — Break curvature-stable regions.

    \item \textbf{Boundary erosion} — Prevent boundary-restored zones from
    forming.

    \item \textbf{Synthetic dominance fields} — Overwhelm coherence-preserving
    manifolds.
\end{itemize}

\medskip

\noindent When Recovery Bands collapse, drift becomes self-propagating. Without
geometric attractors to absorb displacement, drift vectors accelerate (G2.5),
gain directionality, and propagate across roles, contexts, containers, and
scales. Collapse of recovery bands is the geometric indicator that synthetic
pressure has exceeded the manifold’s restorative capacity.

\medskip

\noindent Drift Recovery Bands are therefore the final component of the drift
vector field. They define the geometric conditions under which drift can be
reversed and meaning can return to lawful structure. G2.6 closes the drift
manifold and provides the foundation for the synthetic-pressure dynamics
formalized in G3--G7.
